Prime-degree CM covers give a stronger simultaneous arithmetic
normalization. The distinction between a local reflex-norm image and
the full arithmetic image is essential here.
For each prime \(p\ge5\), put \(h=(p-1)/2\),
\(\mathcal O_p=\mathbb Z[\zeta_p]\otimes\mathbb Z_2\), choose a primitive
residue \(g\), and let \(s=[g]\), \(c=s^h\). Define
\[
 T_p=\{u\in1+2\mathcal O_p:u\,c(u)=1\},\qquad
 P_p=(s-1)T_p.
\]
The algebra \(\mathcal O_p\) includes every factor at two.
The compact-homology scalar of \(y^p=x(1-x)\), restricted to
\(R_{\rm quad}\), will be written \(h_p\); it is the scalar
\(\mathfrak h_p\) used in the earlier cyclic-cover construction.
Define \(Q_p\) by
\[
 \Phi_p^*=\sum_{a=1}^h[a]^{-1},\qquad
 (s-1)\Phi_p^*=(c-1)Q_p,\qquad
 Q_p\in\mathbb Z[s]/(s^h+1).
\]
Its augmentation is \(1\) or \(-1\). Let
\(M_p=(\mathcal O_p/2\mathcal O_p)^c\), with cyclic basis
\(\beta_k=\overline\zeta_p^{\,g^k}+\overline\zeta_p^{-g^k}\),
\(0\le k<h\).

The first condition on \(p\) is that the natural map
\[
 T_p/Q_pT_p\longrightarrow M_p/\overline Q_pM_p,\qquad
 u\longmapsto(u-1)/2\bmod2,
\]
is an isomorphism. Under this condition inclusion identifies
\(P_p/Q_pP_p\) with the same finite elementary abelian group
\(\mathscr C_p\). Write \(\pi_p:M_p\to\mathscr C_p\), and let
\(\mathscr E_p\) be the full unit group of
\(\mathbb Q(\zeta_p+\zeta_p^{-1})\).
The remaining condition is injectivity of
\[
 \mathscr K_p:\mathscr C_p^\vee
       \longrightarrow\mathscr E_p/(\{\pm1\}\mathscr E_p^2),
 \qquad
 \ell\longmapsto
 \left[\prod_{k=0}^{h-1}
   (\zeta_p^{g^k}+\zeta_p^{-g^k})^{\ell(\pi_p\beta_k)}\right].
\]
Let \(\Pi_{\rm global}\) be the set of primes satisfying both
conditions. Square roots in the target may be arbitrary real units.

\begin{theorem}\label{thm:main-cyclotomic-cm-normalization}
For every prime satisfying the first-layer condition, the exact
arithmetic image of \(h_p\) on \(R_{\rm quad}\cap H\) or
\(V_{\rm quad}\cap H\) is the preimage of
\((\ker\mathscr K_p)^\perp\) in \(P_p\). This remains true after
imposing any collection of the other-prime scalar kernels.
For \(p\in\Pi_{\rm global}\), the joint images of
\(R_{\rm quad}\), \(V_{\rm quad}\), and both arithmetic
intersections are all \(\prod_pP_p\), with all finite signs included.
Put
\[
 R_{\rm global}=R_{\rm quad}\cap
          \bigcap_{p\in\Pi_{\rm global}}\ker h_p,\qquad
 V_{\rm global}=V_{\rm quad}\cap R_{\rm global}.
\]
Then
\[
 \begin{gathered}
 G=R_{\rm global}H,\qquad
 R_{\rm global}/V_{\rm global}\simeq R_{\rm quad}/V_{\rm quad},\\
 B\simeq R_{\rm global}/
       \overline{[G,R_{\rm global}](R_{\rm global}\cap H)},\\
 B_{\max2}\simeq
 \left(V_{\rm global}/
       \overline{[G,V_{\rm global}](V_{\rm global}\cap H)}\right)_{\max2}.
 \end{gathered}
\]
The quotient in the first line is pro-odd, and its arithmetic image
is preserved. Both original quadratic restriction kernels are exactly
\(\mathcal S\), with image every invariant quadratic character
vanishing on the arithmetic intersection.

The set \(\Pi_{\rm global}\) contains every prime with odd cyclotomic
class number, and also \(29\) and \(113\); it does not contain \(163\).
At \(29\) and \(113\) the local index is eight but the full arithmetic
image is \(P_p\). At \(163\) the arithmetic image is exactly
\(Q_{163}P_{163}\), of index four.
\end{theorem}

The determinant test is the classical relative class-number formula
in an integral CM-type normalization. The global criterion uses
quadratic Kummer theory of real units and Otsubo's character-labelled
Jacobi-sum comparison. Frobenius at \(59\) and \(227\) fills the local
deficits at \(29\) and \(113\); at \(163\), an explicit negative-square
identity proves that the deficit persists arithmetically.
These arguments, with every finite sign factor, appear in
Sections~\ref{sec:cm-type-unit-criterion}--\ref{sec:simultaneous-global-cm}.
Every added scalar kernel has infinite index.

There is also a restriction statement after all prime scalar kernels,
even when their full arithmetic images are not known to coincide.
For a closed normal subgroup \(L\subseteq N\), put
\[
 B(L)=L/\overline{[G,L](L\cap H)}.
\]

\begin{theorem}\label{thm:main-all-cm-restriction}
Let
\[
 R_{\rm all}=R_{\rm quad}\cap\bigcap_{p\ge5}\ker h_p,\qquad
 V_{\rm all}=V_{\rm quad}\cap R_{\rm all}.
\]
Restriction from \(X_2\) to either group has kernel exactly
\(\mathcal S\). There are canonical maps
\[
 B(V_{\rm all})_{\max2}\simeq B(R_{\rm all})_{\max2}
                         \twoheadrightarrow B_{\max2}.
\]
Moreover \(R_{\rm all}/V_{\rm all}\simeq
R_{\rm quad}/V_{\rm quad}\), preserving the arithmetic image.
The last arrow is asserted to be surjective.
\end{theorem}

Thus every original relative quadratic character remains detectable
on the all-prime kernel. Identifying its intrinsic relative quotient
with \(B_{\max2}\) requires extending those intrinsic characters.
The next theorem gives the exact quadratic extension condition.

\begin{theorem}\label{thm:main-all-cm-central-obstruction}
Let \(\mathcal I\) be any set of primes at least five and put
\[
 R_0=R_{\rm quad},\qquad
 C_p=h_p(R_0),\qquad D_p=h_p(R_0\cap H),\qquad
 J=R_0\cap\bigcap_{p\in\mathcal I}\ker h_p.
\]
The joint full and arithmetic images are exactly
\[
 C=\prod_{p\in\mathcal I}C_p,\qquad
 D_0=\prod_{p\in\mathcal I}D_p.
\]
For the discrete conjugation module
\(U=\operatorname{Hom}_{\rm cont}(C/D_0,\mathbb F_2)\),
\[
 H^j_{\rm cont}(G/R_0,U)=0\qquad(j\ge0).
\]
Let \(\mathscr X(J)\) be the space of continuous \(G\)-invariant
quadratic characters of \(J\) that vanish on \(J\cap H\).
Restriction identifies
\[
 \operatorname{Hom}_{\rm cont}(B,\mathbb F_2)
 \simeq\ker\!\left(
 o_2:\mathscr X(J)\longrightarrow H^2_{\rm cont}(C,D_0;\mathbb F_2)
 \right).
\]
Here \(o_2\) is the relative central extension class, with the
splitting over \(D_0\) specified by the arithmetic subgroup.
Every character in this kernel has a unique extension to \(G\)
which vanishes on \(H\).
\end{theorem}

The proof is Proposition~\ref{prop:all-cm-exact-products},
Corollary~\ref{cor:all-cm-conjugation-vanishing}, and
Theorem~\ref{thm:all-cm-central-extension}, together with the
relative quotient description of \(B\).
The statement includes the infinite collection of all prime degrees.
It does not require a first-layer test or determine the individual
full images \(C_p\). The unknown relative central classes are now
the sole obstruction to extending intrinsic quadratic characters.

At degree \(163\), the possible central obstruction has a particularly
small target.

\begin{theorem}\label{thm:main-163-odd-inertia}
Let \(R=R_{\rm global}\), \(P=P_{163}\), \(A=Q_{163}P\), and
\(U=\operatorname{Hom}_{\rm cts}(P/A,\mathbb F_2)\), with its
cyclotomic conjugation action. Then
\[
 H^j_{\rm cont}(G/R,U)=0\qquad(j\ge0).
\]
If \(h_{163}(R)=P\) and \(J=\ker(h_{163}:R\to P)\), a continuous
\(G\)-invariant quadratic character of \(J\) that is zero on
\(J\cap H\) extends to \(G\), zero on \(H\), exactly when its
relative central extension class is zero. The extension is unique.
The possible invariant relative class is the one-dimensional
quaternion line in \(H^2(P,A;\mathbb F_2)\).
\end{theorem}

The proof is Theorem~\ref{thm:163-concrete-odd-inertia} and
Corollary~\ref{cor:163-single-extension-obstruction}.
The full scalar image is either its arithmetic image \(A\) or \(P\).
The first case gives a further arithmetic normalization. The second
gives a quotient \(A_4\) whose arithmetic image has order three.
The theorem eliminates the conjugation obstruction in the latter
case; the relative central class remains to be determined.
A fixed finite-word detector expresses the full-image alternative.
Neither alternative is selected here, and
\(\operatorname{Hom}(B,\mathbb F_2)\) remains undetermined.

